\section{Introducción: por qué este episodio es una clase sobre la genealogía técnica de los Agent}

Este episodio no es una entrevista común con un invitado, sino una historia técnica de los Agent contada de viva voz. El invitado, Su Yu (苏煜), es profesor del Departamento de Ciencias de la Computación de la Universidad Estatal de Ohio y fundador de NeoCognition; desde hace años investiga temas como Language Agent, Computer Use Agent, Mind2Web, LM Planner y MMMU. La entrevista intenta responder una gran pregunta: por qué en 2026 el relato sobre la IA pasa del Chat al Agent, por qué un producto como OpenClaw provocó en la industria una sacudida parecida al ChatGPT Moment, y por qué se están reorganizando las fronteras entre browser, desktop, mobile, GUI, CLI, API y coding.

\begin{importantbox}{La pregunta central de este episodio}
Si el ChatGPT Moment se entiende como la confirmación pública del paradigma de los LLM, el OpenClaw Moment se parece más a la confirmación pública del paradigma de los Agent: el modelo ya no solo responde preguntas, sino que entra en un entorno, usa herramientas, ejecuta acciones, aprende de forma continua y se acerca poco a poco a un universal digital agent.
\end{importantbox}




\subsection{Resumen de la sección}

Este episodio se presta para organizarse como una nota sobre la «genealogía técnica de los Agent»: primero se define qué es un Agent; después se repasan el logical agent, el neural agent, el semantic parsing y el language agent; por último se discuten OpenClaw, el universal digital agent, el continual learning, el world model y la difusión en la industria.
